\documentclass[11pt]{article}
\usepackage{amsmath,amssymb,amsthm,fullpage,hyperref}
\newcommand{\HH}{\mathbb{H}}
\newcommand{\CC}{\mathbb{C}}
\newcommand{\ZZ}{\mathbb{Z}}
\newcommand{\SLZ}{{\rm SL}_2(\ZZ)}
\newcommand{\tk}{\widetilde{k}_T}
\newcommand{\Li}{{\rm Li}}
\newcommand{\code}[1]{\texttt{#1}}

\title{Numerical validation of \code{finite\_contour\_cocycles\_short.tex}}
\author{Companion to \code{validate.py}}
\date{}

\begin{document}
\maketitle

\section{What this suite is for}

Every closed form in the note is a claim that some explicitly-summable expression equals a
contour integral. Algebra alone does not catch a claim that is right in one geometric regime
and wrong in another: the error is invisible until the wrong regime is evaluated. This suite
therefore takes \emph{raw quadrature of the defining contour integral} as ground truth and
checks each closed form against it, never against another closed form.

Two design rules follow.

\paragraph{Index by regime, not by equation.} Each check runs across the geometric
configurations the note distinguishes: a pole above the $T$-segment, on the edge
${\rm Im}\,z=1$ with $z\neq i$, at $i$, below the contour, and absent. A per-equation test that
samples one configuration certifies nothing about the others.

\paragraph{Ground truth must be exact.} Where a check needs Fourier coefficients, they come from
integer $q$-series arithmetic, not from numerical extraction. Recovering $c_f(n)$ from
$\int_0^1 f(x+iy)e^{-2\pi in(x+iy)}dx$ is catastrophically ill-conditioned for a form with a pole
at the cusp: for $1/\Delta$ the integrand runs as $e^{2\pi ny}$ while the answer is only
$e^{4\pi\sqrt n}$, so $c(18)$ at $y=1.3$ asks for a $10^{23}$ number out of $10^{67}$-sized
values. The suite computes $\prod(1-q^n)^{\pm24}$ in exact integer arithmetic instead;
\code{cf\_Delta} and \code{cf\_invDelta} reproduce $\tau(n)$ and the known
$1/\Delta=q^{-1}+24+324q+3200q^2+\dots$ exactly.

Each check carries the \verb|\label| it validates, so a failure names a line in the note rather
than a formula in Python. Two tiers: \code{validate.py} (coarse precision, one case per regime)
and \code{validate.py --slow} (high precision, full cross-product). A third tier,
\code{report}, prints a measured discrepancy without passing or failing; it is used for
questions the note has not yet settled, so the suite documents them instead of presuming an
answer.

\subsection*{Shared machinery in \code{validate.py}}

\noindent\begin{tabular}{@{}ll@{}}
\code{ktil(tau,s,k)} & the Hurwitz $T$-kernel $\tk$ of \code{def:kernel}\\
\code{LT\_raw(f,k,s,tau0)} & \textbf{ground truth}: quadrature of
   $e^{-i\pi s/2}\!\int_{\tau_0-1}^{\tau_0}\! f\,\tk\,d\tau$\\
\code{LT\_closed(cf,k,s,tau0)} & the closed form \code{eq:T-segment-closed-Mk}\\
\code{E4}, \code{E6}, \code{Delta}, \code{invDelta} & the forms as functions on $\HH$, for quadrature\\
\code{cf\_E4}, \code{cf\_E6}, \code{cf\_Delta} & their \emph{exact} integer $q$-coefficients\\
\code{cf\_invDelta}, \code{cf\_j}, \code{jay} & likewise for $1/\Delta$ and $j=E_4^3/\Delta$\\
\code{residue\_at(g,s0,r)} & $\mathop{\rm Res}$ by trapezoid rule on a circle $|s-s_0|=r$\\
\code{WEAK\_FORMS} & the Layer C form table (form, function, coefficients, $k$)\\
\code{\_eta\_prod}, \code{\_smul}, \code{\_sinv} & integer power-series arithmetic behind those\\
\code{check(label,layer,tier)} & decorator registering a test against a \verb|\label|\\
\code{rel(got,want)} & relative error, denominator floored at $1$\\
\code{main} & driver; \code{TOL}, \code{NT}, \code{NMAX} are set by tier\\
\end{tabular}

\section{Layer A: kernel and unfolding}

Layer A sits upstream of everything else. If the kernel identities or the mode-by-mode
unfolding are wrong, every downstream check is measuring the wrong thing.

\paragraph{A1 --- the $\delta_T$ shift (\code{def:kernel}; \code{a1\_kernel\_shift}).}
Verifies
$\delta_T\tk(\tau,s)=\tk(\tau+1,s)-\tk(\tau,s)=-(\tau+1)^{s-1}+e^{i\pi(s-1)}(\tau+1)^{k-1-s}$
directly, at three weights, three base-points and three values of $s$ including one on the
critical line. This is the identity that turns the $\partial_{\tau_0}$ boundary term into a unit
shift of the kernel, so it underwrites $\tau_0$-independence.

\paragraph{A2 --- the Hurwitz antiderivative (\code{eq:hurwitz-antideriv}; \code{a2\_hurwitz}).}
Three separate claims: $\partial_a\big[\zeta(\sigma-1,a)/(1-\sigma)\big]=\zeta(\sigma,a)$; the
shift $\zeta(\sigma-1,a+1)-\zeta(\sigma-1,a)=-a^{1-\sigma}$; and the definite integral
$\int_{i-1}^{i}\zeta(\sigma,\tau+1)d\tau=i^{1-\sigma}/(\sigma-1)$ that they combine to give.
These produce the $n=0$ contribution, hence the entire
$-c_f(0)\big(\tfrac1s+\tfrac{i^k}{k-s}\big)$ polar structure.

\paragraph{A3 --- the unfolding (\code{eq:T-segment-closed-Mk}).}
The general-$\tau_0$ $T$-segment closed form, checked against quadrature in three stages of
increasing exposure:
\begin{itemize}
\item \code{a3a\_unfold\_pos} --- positive modes only ($\Delta$, $k=12$): no branch question arises;
\item \code{a3b\_unfold\_const} --- $c_f(0)\neq0$ ($E_4$, $k=4$): exercises the
      $-c_f(0)\big((\tau_0/i)^s/s+\dots\big)$ term;
\item \code{a3c\_unfold\_neg} --- negative modes ($1/\Delta$, $k=-12$): $n<0$ makes $(2\pi n)^s$ a power of a negative
number and drives $-2\pi in\tau_0$ onto the cut of $\Gamma(s,\cdot)$.
\end{itemize}

The third stage is the ``incomplete-gamma unfolding ambiguity''. The suite pins the convention
by quadrature: for $n<0$,
\[
(2\pi n)^s := (2\pi|n|)^s\,e^{+i\pi s}.
\]
This must be applied as an explicit phase and never folded into the base as
$(2\pi|n|e^{i\pi})^s$: numerically $e^{i\pi}=-1+\varepsilon i$ with the sign of $\varepsilon$ a
rounding artifact, so the base-folded form silently changes branch with the working precision.
With the phase explicit, the closed form matches quadrature to $10^{-31}$ at
$\tau_0=i,\,1.35i,\,2i$ and at $\tau_0$ off-axis to the right.

\paragraph{A3, ${\rm Re}\,\tau_0<0$: the continued $\Gamma$ (\code{a3d\_unfold\_neg\_left}).}
For $n<0$ write $z=-2\pi in\tau_0$. With $\tau_0=x+iy$ this is
$z=-2\pi|n|y+2\pi i|n|x$, so ${\rm Re}\,z<0$ always while ${\rm Im}\,z$ carries the sign of
$x={\rm Re}\,\tau_0$: the point $z$ sits just above the negative-real cut of $\Gamma(a,\cdot)$
when ${\rm Re}\,\tau_0>0$ and just below it when ${\rm Re}\,\tau_0<0$. The L-integral varies
continuously as $\tau_0$ moves; the \emph{principal} $\Gamma$ jumps.

The jump is the standard monodromy. Since $\Gamma(a,z)=\Gamma(a)-\gamma(a,z)$ with
$\gamma(a,z)=z^a\sum_{m\ge0}(-z)^m/\big(m!(a+m)\big)$, all the multivaluedness sits in $z^a$,
and one anticlockwise turn gives
\[
\Gamma_{\rm cont}(a,z)=\Gamma(a)\big(1-e^{2\pi ia}\big)+e^{2\pi ia}\,\Gamma_{\rm prin}(a,z),
\]
implemented as \code{\_gamma\_cont} and applied per term at $a=s$ and $a=k-s$. With it,
\code{eq:T-segment-closed-Mk} reproduces quadrature to $10^{-31}$ for ${\rm Re}\,\tau_0<0$,
just as the principal branch does for ${\rm Re}\,\tau_0\ge0$.

So the lemma is correct as stated for all $\tau_0\in\HH$, and \code{prop:HGG}'s phrase ``the
identity continues to $n<0$'' is doing real work: $\Gamma$ must be the continuation, which
coincides with the principal value only while ${\rm Re}\,\tau_0\ge0$. An earlier draft of this
document reported the lemma as failing for ${\rm Re}\,\tau_0<0$; that was an artifact of
evaluating $\Gamma$ on the principal branch on the far side of its cut, not a defect of the
result.

\section{Layer B: $\tau_0$-independence and wall-crossing}

Layer B tests the full L-integral $L^*=L^*_S+L^*_T$ of \code{def:Lint}, with both segments by
quadrature (\code{LS\_raw}, \code{LT\_raw\_path}, assembled by \code{L\_raw}). Paths are
polylines integrated by \code{path\_int}; $\HH$ is convex, so the straight line from $-1/\tau_0$
to $\tau_0$ stays inside it.

\paragraph{B1, B2 --- $\tau_0$-independence (\code{b1\_indep\_onaxis}, \code{b2\_indep\_offaxis}).}
\code{prop:indep}: neither segment is separately invariant, only the sum. Checked on the
imaginary axis and off it, for $\Delta$, $E_4$, $E_6$ and $1/\Delta$. B2 includes
${\rm Re}\,\tau_0<0$ deliberately, which is what first showed that the A3 discrepancy could not
be a property of the L-integral: the raw $L^*$ is invariant there to $10^{-16}$, so whatever was
wrong had to lie in the evaluation of the closed form --- as it did.

\paragraph{B3 --- a meromorphic form (\code{b3\_indep\_meromorphic}).}
\code{prop:indepHomotopy} for $f=\Delta/(j-j(2i))\in F_{12}$, whose poles are the $\Gamma$-orbit
of $2i$ at heights $2,\tfrac12,\tfrac25,\dots$. Base-points near $i$ keep both segments clear of
every image, so these all lie in one class and $L^*$ must not move.

\paragraph{B4, B5 --- wall-crossing (\code{b4\_wall\_T}, \code{b5\_wall\_S}).}
\code{lem:wall} with \emph{fixed end-points and two homotopy classes}, so the crossing is
isolated from the genuine $\tau_0$-dependence of each separate segment. For $T$: a straight
segment at height $2.2$ against one dipping to $1.8$, enclosing $2i$ alone. For $S$: two paths
agreeing around every image except $2i$, which one passes to the left and the other to the
right. In both the enclosed loop runs clockwise, so $X=-1$ and the predicted jump is
$-2\pi ie^{-i\pi s/2}r$, with $r_T={\rm Res}(f\tk)$ and $r_S={\rm Res}(f\tau^{s-1})$ from
\code{res\_circle}.

Two things had to be got right for these to test anything. The form is normalised so
${\rm Res}_{2i}f=1$; unnormalised the residue is $\Delta(2i)/j'(2i)\approx2\times10^{-12}$ and
every jump is $\sim10^{-9}$, at which scale a sign error hides under the noise floor. And
\code{res\_circle} needs enough points: the trapezoid rule on $|\tau-2i|=r$ converges like
$(r/R)^M$ with $R=1$ the distance to $2i+1$, so the original $r=0.3$, $M=16$ gave only
$0.3^{16}\sim4\times10^{-9}$ --- agreement to eight digits, which reads as a failure against a
$10^{-9}$ tolerance and invites a hunt for a nonexistent bug elsewhere.

\section{Layer C: \code{thm:weakL} and the functional equation}

For $f\in M^!_k$ the base-point may be taken at $\tau_0=i$, where the $S$-segment degenerates to
a point and $L^*=L^*_T$ alone. Ground truth for this layer is therefore
\code{LT\_raw(f,k,s,I)}, the bare quadrature of $e^{-i\pi s/2}\!\int_{i-1}^{i}f\,\tk\,d\tau$.

The five test forms (\code{WEAK\_FORMS}) are chosen to span the cases the statement
distinguishes rather than to be interesting individually:
\begin{center}
\begin{tabular}{@{}llll@{}}
form & code & $k$ & feature exercised\\
\hline
$\Delta$    & \code{cf\_Delta}    & $12$  & cusp form, $c_f(0)=0$\\
$E_4$       & \code{cf\_E4}       & $4$   & $c_f(0)\ne0$, so the polar term is live\\
$E_6$       & \code{cf\_E6}       & $6$   & $i^k=-1$\\
$1/\Delta$  & \code{cf\_invDelta} & $-12$ & pole at the cusp, so $n<0$ modes\\
$j$         & \code{cf\_j}        & $0$   & the $s=0$ and $s=k$ poles coincide\\
\end{tabular}
\end{center}
\noindent$j=E_4^3/\Delta$ is built by integer series arithmetic and reproduces
$q^{-1}+744+196884q+21493760q^2$ exactly; $E_6$ reproduces $1-504q-16632q^2-122976q^3$.

\paragraph{C1 --- \code{eq:weakL} (\code{c1\_weakL}).}
The closed form against quadrature, for all five forms. Since \code{eq:weakL} is
\code{eq:T-segment-closed-Mk} specialised to $\tau_0=i$, this also re-exercises the branch
convention pinned in A3, now with $c_f(0)\neq0$ and with $k\le0$.

\paragraph{C2 --- the residue (\code{c2\_residue}).}
$\mathop{\rm Res}_{s=0}L^*=-c_f(0)$. The residue is extracted from the \emph{raw} integral by a
contour in $s$ (\code{residue\_at}, trapezoid rule on $|s|=r$), so the note's polar claim is
tested against quadrature and not against its own closed form. The circle is used rather than
extrapolating $s\,L^*(s)$, which converges only linearly in $s$ and would not reach tolerance.

The weight-zero case is the sharp one. For $k=0$ the two poles collide, and the polar term is
$-c_f(0)\big(\tfrac1s+\tfrac{i^0}{0-s}\big)=0$ identically, so the residue must \emph{vanish}
even though $c_j(0)=744$. It does. This is the mechanism behind the factor $(1-\delta_{k,0})$
appearing in the special-value corollaries of arXiv:1806.09874, recovered here from bare
quadrature.

\paragraph{C3 --- the functional equation (\code{c3\_functional\_eq}).}
$L^*(f,s)=i^kL^*(f,k-s)$, with \emph{both} sides computed by raw quadrature. This is the
cheapest global probe in the suite: it touches every phase in the construction at once, and a
single lost sign anywhere breaks it. $E_6$ is the case that matters, being the only form here
with $i^k=-1$; the weights $12,4,-12,0$ all give $i^k=+1$ and would not detect a dropped
$i^k$.

\section{Layer D: periods and quasi-periods}

Two of the three checks here are \emph{external}: they compare the note's construction against
classical objects computed by an independent route, rather than against another of the note's
own formulas. The test forms are the unique normalised cusp forms $\Delta_k=q+O(q^2)$ at the
weights with $\dim S_k=1$ (\code{CUSP1}: $k=12,16,18,20,22,26$), built by exact integer series
multiplication from $\Delta$, $E_4$, $E_6$; $\Delta_{16}=\Delta E_4$ reproduces
$q+216q^2-3348q^3+\dots$ exactly.

\paragraph{D1 --- $r_f\in W$ (\code{d1\_rfW}).}
\code{def:rf} assembles the period polynomial from the $L^*$-values,
$r_f=(2\pi i)^{n+1}\sum_\ell(-1)^\ell\binom{n}{\ell}i^{\ell+1}L^*(f,\ell+1)X^{n-\ell}Y^\ell$, and
\code{lem:rfW} claims it lies in $W$. The test applies the slash action
$(w|\gamma)(X,Y)=w(aX+bY,cX+dY)$ (\code{slash}) at $S$ and $U=TS$, and checks
$r_f|_{(1+S)}$ and $r_f|_{(1+U+U^2)}$ vanish relative to the size of $r_f$. They do: the
$S$-relation to machine zero, the $U$-relation to $10^{-11}$--$10^{-13}$. So the period
polynomial built from the finite two-segment contour satisfies the Eichler--Shimura relations.

\paragraph{D2 --- the classical periods (\code{d2\_classical\_periods}).}
\textbf{External anchor.} $r_m=i^{m+1}L^*(m{+}1)$ must be the classical critical period
$\int_0^{i\infty}\tau^mf(\tau)\,d\tau=i^{m+1}\int_0^\infty y^mf(iy)\,dy$, i.e.\ $L^*$ must
\emph{be} the completed period $\Lambda(f,s)$. The right side is obtained by direct quadrature
on the axis, not by re-deriving the incomplete-gamma sum, so this fixes the normalisation of
\code{def:rf} independently. Agreement is $10^{-11}$--$10^{-16}$ at both ends of the critical
strip.

The integrand must be evaluated through \code{f\_on\_axis}. A cusp form decays at both ends of
the axis, but the $q$-series does not converge at the lower one: at $y=0.03$ one has
$|q|=0.83$, where forty terms of $E_4$ are meaningless. Below $y=1$ the evaluation is pulled
back by $f(iy)=f(i/y)/(i^ky^k)$, which lands it above the unit circle. Without that the
computed period comes out as $6\times10^{11}$ in place of $10^{-73}$.

\paragraph{D3 --- Haberland--Petersson (\code{d3\_haberland}).}
\textbf{The strongest check in the suite}, because it is normalisation-\emph{invariant} and
both sides are computed here from scratch. Cohen, ANTS X (2013), Thm 5.2(2), for the full
modular group:
\[
-6(-2i)^{k-2}\langle f,f\rangle=\sum_{m+n\le k-2}\binom{k-2}{m+n}\binom{m+n}{m}(-1)^m
{\rm Im}\big(r_m\overline{r_n}\big),
\]
with $r_m=i^{m+1}L^*(f,m{+}1)$. The left side is the Petersson norm by direct integration of
$|f|^2y^{k-2}$ over the fundamental domain; the right uses only the finite-contour $L^*$. They
agree to $6\times10^{-16}$.

Unlike the determinant ratios of Table~1, which depend on how the period-polynomial basis is
normalised, nothing here can be rescued or spoiled by a choice of basis: the pairing is
contracted away on both sides. That makes it the closest thing in this suite to an outside
verdict on the construction.


\section{Layer E: the meromorphic machinery of \S4}

This is the layer with the most apparatus and the least external corroboration, so it carries
two ground truths of its own. \code{IN\_raw} evaluates $I_N(s,z)$ directly: by
\code{prop:indepHomotopy} the $\delta\to0$ representative of the minimal class may be used, and
there the block is regular at $i$ (as ${\rm Im}\,z>1$), so the $S$-integral vanishes and the
$T$-segment is $[i-1,i]$. \code{wallvalue\_raw} evaluates $L^*(\widehat P_if,s;i(1+\delta))$ at
finite $\delta$, the $S$-segment split into a regular part and the principal value of the pure
pole, as \code{def:null\_homotopy} prescribes. Three meromorphic forms with a pole at $i$ span
the pole orders (\code{MERO\_AT\_I}): $\Delta/E_6$ ($k=6$, $P=1$), $E_4^2\Delta/E_6^2$ ($k=8$,
$P=2$), $\Delta^2/E_6^3$ ($k=6$, $P=3$) --- orders and weights consistent with
\code{prop:ellord}.

\paragraph{E1 --- the $\Gamma$ recursion (\code{e1\_gammarec}).}
\code{eq:gammarec} at three cutoffs $J$, three arguments $X$ (real, complex, large) and two
$s$. Cheap, and upstream of both \code{lem:polylogedge} and \code{lem:wallvalue}.

\paragraph{E2 --- the block expansion (\code{e2\_blockexp}).}
\code{eq:blockexp} claims the block is a pure order-$m$ pole plus an entire $T$-periodic tail.
The test forms the remainder $\tfrac{(-2\pi i)^m}{(m-1)!}\Li_{1-m}(e^{2\pi i(\tau-i)})-(\tau-i)^{-m}$
and checks its residue at $i$ vanishes, for $m=1,2,3$. That vanishing is precisely what makes
$r^*_{f,z}(m)$ the Laurent coefficients rather than merely proportional to them, so it validates
the Option-A reading of \code{def:proj}.

\paragraph{E3 --- $I_N$ above the line (\code{e3\_polylogcf}).}
\code{lem:polylogcf} against \code{IN\_raw}, at $N=0,1,2$ and $z$ both on and off the imaginary
axis. Everything in \S4 is built on this lemma, and $N=0$ is the case carrying the
$\delta_{N,0}\big(\tfrac1s+\tfrac{i^k}{k-s}\big)$ inversion constant --- the term responsible
for the $1/s$ poles that distinguish this $L^*$ from the cusp-anchored one of
arXiv:1806.09874. It reproduces quadrature to $10^{-16}$.

\paragraph{E4 --- the wall value (\code{e4\_wallvalue}).}
\code{lem:wallvalue} asserts a \emph{limit}, so the raw integral must be extrapolated rather
than sampled: $L=\big(8f(\delta/4)-6f(\delta/2)+f(\delta)\big)/3$ annihilates both an
$O(\delta)$ and an $O(\delta^2)$ error term. Both orders occur, and which one appears is not
incidental --- it splits along the parity dichotomy of \code{prop:ellord}:
\begin{center}
\begin{tabular}{@{}lll@{}}
form & regime & convergence\\
\hline
$E_4^2\Delta/E_6^2$ & $4\mid k$, $P$ even & $O(\delta^2)$, ratio $\to4$ per halving\\
$\Delta/E_6$, $\Delta^2/E_6^3$ & $k\equiv2\bmod4$, $P$ odd & $O(\delta)$, ratio $\to2$\\
\end{tabular}
\end{center}
\noindent All three converge to $\sum_m\tfrac{(-2\pi i)^m}{(m-1)!}r^*_{f,i}(m)\mathrm{FP}_{m-1}(s,k)$.
The tolerance here is $10^{-5}$, not the global one: what survives extrapolation is the
$O(\delta^3)$ term the three-point formula cannot remove. The substantive evidence is the
convergence \emph{rate}, not the residual.

\paragraph{E5 --- the cutoff (\code{e5\_cutoff}, \code{e5b\_cutoff\_minimal}).}
\code{rem:cutoff} holds that $\mathrm{FP}_N$ is independent of $J$ for $J\ge N+1$. It is, but
each $J$ carries its own truncation: the remainder decays as $n^{-(J-N+1)}$, so at a fixed
$n_{\max}$ only the larger cutoffs are numerically converged. The equality is asserted among
those; the slower ones, including the minimal $J=N+1$, are shown converging in a
\code{report}-tier table --- at $J=N+1$ the error falls as $1/n_{\max}$, tripling $n_{\max}$
thirding the error, exactly the $n^{-2}$ tail the note predicts. \code{FP} therefore
\emph{evaluates} at $J=N+5$, which is legitimate precisely because the lemma being tested is
true.

\paragraph{E6 --- the modularity relations (\code{e6\_laurent}).}
\code{eq:Laurent} at every order $r=1,\dots,P$ for all three forms; holds to machine zero.

\paragraph{E7 --- the edge route (\code{e7\_edge\_route}, \code{report} tier).}
A pole at ${\rm Im}\,z=1$ with $z\neq i$ admits two readings. \code{lem:polylogedge} continues
the from-\emph{above} formula of \code{lem:polylogcf} down to the edge, which is the convention
the note adopts and retains. The from-\emph{below} value --- the block's own $q$-series against
$G_{s,k}$, with no inversion and hence no $\delta_{N,0}$ term --- differs from it by the
crossing residue, measured here as a factor of about $4.4$ at $z=0.3+i$. The line is printed
rather than scored: it records the size of a convention choice, not a defect.

\paragraph{E8 --- the polar defect (\code{e8\_polar\_defect}).}
\code{lem:rfWFk} and \code{def:Qf}: for meromorphic $f$ with simple poles inside the horocyclic
triangle $\mathcal T$ and vanishing strip-constant, $r_f|_{(1+S)}=0$ while
$r_f|_{(1+U+U^2)}=-Q_f$, the residue sum over the enclosed poles. The test form is
$f=\Delta^2/(E_4^3-j(z_0)\Delta)$ with $z_0=0.5+0.8i$, whose enclosed set is the single
$U$-orbit $\{z_0,\,Uz_0,\,U^2z_0\}$ --- the loop is $U$-symmetric, so partial orbits cannot be
enclosed, and the third member $U^2z_0=0.562+0.899i$ is the pole a translate-window orbit
search once missed (a $15\times$ discrepancy until it was found). The identity was pinned at
$\mathrm{dps}=25$ to $10^{-24}$ three independent ways, including direct quadrature of the
loop via the exact leg parameterisations $\tau=t{+}1$, $t/(t{+}1)$, $-1/t$ with $t$ on the
$T$-segment (\code{resolve13.py}); the suite carries the fast-tier version, plus a
non-vacuousness guard that the defect is $O(1)$ before cancellation.

\paragraph{E9 --- the explicit corrector (\code{e9\_explicit\_corrector}).}
\code{thm:rtildeW}: the corrected period polynomial
\[
\tilde r_f=r_f-c_f\,r_{E_k}
+\sum_{\mathcal O}2\pi i\,r^*_{f,q_{\mathcal O}}(1)
\big[(2\pi i)^{n+1}\mathbf{K}_T(q_{\mathcal O};X,Y)-r_{E_k}(X,Y)\big]
\]
lies in $W$: winding $\hat\gamma^T_f$ once about $q_{\mathcal O}=Sp_{\mathcal O}$ (one pole per
enclosed $U$-orbit) trivialises the $(TS)^3$-loop --- the added loop transports under $S^{-1}$
to $+1$ windings on the whole orbit, cancelling the clockwise boundary passage --- at the
explicit price of \code{lem:wall}, and $r_{E_k}$ absorbs the shifted strip-constant. Checked
in both regimes: $f_{z_0}$ (enclosed orbit, $c_f=0$; also with a second choice of
representative, which shifts $\tilde r_f$ within $W$) and $f_7$ ($c_f\ne0$, nothing enclosed,
correction $=-c_fr_{E_k}$). At $\mathrm{dps}=25$ both Eichler--Shimura relations hold to
$10^{-22}$ (\code{explicitW.py}); the suite carries the fast-tier version. This is the
statement toward which the whole layer builds: period polynomials in $W$ for arbitrary
simple-pole $f\in F_k$, from residue data alone.

\paragraph{A note on tolerances.} Two Layer E checks carry their own tolerance instead of the
global \code{TOL}. This is deliberate. \code{TOL} scales with the working precision, but a
check limited by truncation or by extrapolation does not improve when mpmath is handed more
digits: at $\mathrm{dps}=30$ the global tolerance is $10^{-24}$, while $\mathcal E^+$ truncated
at $n_{\max}=60$ caps agreement near $10^{-15}$ however many digits are carried. A tolerance
must reflect the numerics of the test that uses it, not the width of the mantissa.

\section{Layer F: external anchors}

Layers A--E check the note against itself and against raw quadrature. Layer F checks it against
things outside it: the earlier paper, and a statement about how the two must differ.

\paragraph{F1 --- the published $y<B$ branch (\code{f1\_1806\_below}).}
\code{eqLemRegX3} of arXiv:1806.09874 in the case where the ray never crosses the circle
$|e(it-\tau_p)|=1$. It is correct as published, reproducing \code{I1806\_raw} to $10^{-16}$ for
$N=0,1,2$. It also pins that paper's conventions: the integrand is the proof's
$\tfrac{dt}{t}t^{s}$, not the definition's $\tfrac{d\tau}{\tau}(\tau/i)^{s-1}$.

\paragraph{F2 --- the corrected $y>B$ branch (\code{f2\_1806\_above}).}
Regression test for the corrigendum in \code{fix\_previous\_1806\_file/}. For $y>B$ the ray
crosses the circle at $t=y$, and neither the $q$-series nor its inversion is valid across that
point, so the integral must be split there. The published version evaluates a single
antiderivative at $t=B$ and $t=\infty$, which drops three things: the $t=y$ end-point, the
$\delta_{N,0}$ constant from the $N=0$ inversion, and the whole $\int_y^\infty$ tail. The
corrected form reproduces quadrature to $10^{-6}$; the published one misses by a factor of
about $12$.

Two implementation notes, both of which cost a wrong answer before being fixed. The branch is
the one the corrigendum specifies by
$\big[\Gamma(s,-2\pi nB)-\Gamma(s,-2\pi ny)\big]/(-2\pi n)^s=\int_B^y\tfrac{dt}{t}t^se^{2\pi nt}$;
with $(-2\pi n)^s$ read principally as $(2\pi n)^se^{+i\pi s}$, dividing contributes
$e^{-i\pi s}$, which matches that integral to $10^{-25}$. Pairing $e^{+i\pi s}$ with mpmath's
default branch instead converges to a value $26\%$ off. And the two boundary series converge
only conditionally --- their terms go as $n^{N-1}$ times a phase --- so partial sums are
Ces\`aro-averaged, which is why the tolerance here is $10^{-4}$ rather than the global one.

\paragraph{F3 --- entirety of the 1806 block integrals (\code{f3\_1806\_entire}).}
Thm~1 of that paper asserts the only non-regular terms are
$-c_f(0)\big(\tfrac{B^s}{s}+\tfrac{i^kB^{k-s}}{k-s}\big)$. That needs the block integrals to be
regular in $s$, which is automatic --- a finite lower end-point plus exponential decay at the
cusp makes $\int_B^\infty$ converge for every $s$ --- and is checked here as a residue in $s$
taken around the raw integral. It vanishes to $10^{-17}$. This is why the defect in F2, which
is a finite-part error, cannot disturb Thm~1, the Hurwitz class-number sum rule, or any of that
paper's special values: all of those are residue statements, and $1/\Gamma(s)$ annihilates
every regular term at $s=0$.

\paragraph{F4 --- the residue--strip rule (\code{f4\_residue\_strip}).}
The sharp form of how this $L^*$ and the cusp-anchored one differ. $L^*_S$ is regular in $s$,
so the pole at $s=0$ comes entirely from the $T$-segment, where the kernel's Hurwitz pole
$\tk\sim-1/s$ multiplies the constant Fourier mode. Hence
\[
\mathop{\rm Res}_{s=0}L^*(f,s)=-\int_{\tau_0-1}^{\tau_0}f(\tau)\,d\tau ,
\]
the integral over one horizontal period being exactly that constant mode --- read in
\emph{the strip the contour sits in}, not at the cusp. For $\Delta/(j-j(2i))$, whose pole at
$2i$ lies above the contour, this holds to $2\times10^{-11}$; for $E_4$, $E_6$ and $1/\Delta$
to $10^{-14}$.

This is what dissolves the apparent conflict with arXiv:1806.09874. That paper anchors its
contour at the cusp, above every pole, so it sees only $c_f(0)$; this one sits at height one
and additionally sees $2\pi i\sum a_{-1}$ from the poles above it. Both are right about their
own object, and the difference is exactly the residues crossed in moving the base-point. A cusp
form is excluded from the check: there both sides vanish identically and a relative comparison
carries no information.

\paragraph{F5 --- the Hurwitz class numbers and traces (\code{f5\_hurwitz\_qexp}).}
\textbf{External arithmetic anchor.} For class number one the Hilbert class polynomial is
linear, so $\Lambda_d=-q\,\partial_q\log\mathcal H_d(j(\tau))$ is elementary:
\[
\Lambda_3=\frac{E_6}{3E_4},\qquad
\Lambda_7=\frac{E_4^2E_6}{E_4^3+3375\,\Delta},
\]
using $q\,\partial_qj=-E_4^2E_6/\Delta$ and $j+3375=(E_4^3+3375\Delta)/\Delta$, each normalised
so that $c(0)=H(d)$. Read off at height $2.2$, above every pole, the $q$-expansion must be
$\Lambda_d=H(d)+\sum_nt_n(d)q^n$ with $H(d)$ the Hurwitz class number and $t_n(d)$ the trace of
$J_n$ over CM points of discriminant $-d$:
\begin{center}
\begin{tabular}{@{}lllll@{}}
& $c(0)$ & $H(d)$ & $c(1)$ & $t_1(d)$\\
\hline
$\Lambda_3$ & $0.333333333333$ & $1/3$ & $-248.0$ & $\big(j(\rho)-744\big)/3=-248$\\
$\Lambda_7$ & $1.0$ & $1$ & $-4119.0$ & $j(\alpha_7)-744=-4119$\\
\end{tabular}
\end{center}
\noindent Three of the four agree to exactly zero. This is Zagier's arithmetic recovered by
contour integration, and it is the one place where the construction is measured against
mathematics that predates it.

\paragraph{F6 --- the Hurwitz sum rule (\code{f6\_hurwitz\_sumrule}).}
Eq.~(s3e2) of arXiv:1806.09874 states $\lim_{s\to0}L^{\rm reg}_{\Lambda_d}(s)=-H(d)$. Since
$L^{\rm reg}=(2\pi)^sL^*/\Gamma(s)$ and $1/\Gamma(s)$ has a \emph{zero} at $s=0$, the limit
picks out the residue alone, so by F4 it equals $-c_f(0)$ read in the strip the contour sits
in. The two discriminants then behave completely differently:
\begin{itemize}
\item $d=3$: the orbit of $\rho$ has maximal height $\sqrt3/2<1$, so every pole lies below the
contour, the strip constant \emph{is} the cusp constant, and the residue is $-H(3)=-1/3$ ---
the sum rule exactly as published.
\item $d=7$: the CM point $\alpha_7=(1+i\sqrt7)/2$ has height $\sqrt7/2>1$ and lies
\emph{above} the contour. A simple zero of $j-j(\alpha)$ always carries
$a_{-1}=-1/(2\pi i)$, so it contributes $2\pi ia_{-1}=-1$ to the strip constant, which cancels
$H(7)=1$ exactly. The residue \textbf{vanishes}, to $1.4\times10^{-16}$.
\end{itemize}
Same form, same machinery, and the $1/s$ residue is $-H(d)$ or $0$ purely according to which
side of the contour the CM point falls on. This is the entire difference between this $L^*$ and
the cusp-anchored one of arXiv:1806.09874, reduced to a single number --- and it is why the
apparent conflict between the two was never a contradiction.

\paragraph{F7 --- lifting the contour above every pole (\code{f7\_raised\_contour}).}
The test behind \code{rem:raised}. Since $L^*_S$ is a finite integral of $f\,\tau^{s-1}$ and is
entire in $s$, the pole at $s=0$ is carried entirely by the $T$-segment, where $\tk\sim-1/s$
multiplies $\int_{\hat\gamma_f^T}f\,d\tau$. Deforming $\hat\gamma_f^T$ to pass above every pole
of $f$ should therefore return the constant term at the \emph{cusp}, recovering the residue
bookkeeping of arXiv:1806.09874. Two cases, chosen so that lifting matters in one and not the
other:
\begin{center}
\begin{tabular}{@{}llll@{}}
& poles vs.\ contour & flat $\mathop{\rm Res}$ & lifted $\mathop{\rm Res}$\\
\hline
$\Lambda_7$ & $\alpha_7$ at height $\sqrt7/2>1$, above & $0$ & $-1=-H(7)$\\
$\Lambda_3$ & orbit of $\rho$ at $\sqrt3/2<1$, all below & $-1/3$ & $-1/3=-H(3)$\\
\end{tabular}
\end{center}
\noindent In both the lifted value is $-c_f(0)$ at the cusp. For $\Lambda_7$ that is a genuine
change of class, and the shift is the crossing of Lemma~\ref{lem:wall}; for $\Lambda_3$ nothing
is crossed and the lift is a homotopy, so the value must not move, which it does not. The
base-points are chosen so the vertical legs of the lifted path miss the poles: the orbits of
$\alpha_7$ and of $\rho$ both sit at ${\rm Re}=\tfrac12\bmod1$, so legs at ${\rm Re}=-1,0$ are
clear.

A form whose constant term vanishes at the cusp is useless for this check. $\Delta/(j-j(2i))$
begins at $q^2$, so both sides are zero and the comparison carries no information; the same
consideration excludes cusp forms from F4.

\section{Layer G: suite hygiene}

The suite is the arbiter, so it is itself checked. Each of these recomputes a quantity with one
numerical knob changed and demands the answer not move.

\paragraph{G1 --- truncation (\code{g1\_truncation}).}
Mode sums under $n_{\max}\to2n_{\max}$, for $L^*$ and for $\mathrm{FP}_N$. This is the failure
that made $\mathrm{FP}$ at the minimal cutoff $J=N+1$ look like a violation of
\code{rem:cutoff}.

\paragraph{G2 --- precision (\code{g2\_precision}).}
Values recomputed at $\mathrm{dps}+12$. This catches results that are artifacts of the mantissa
rather than of the mathematics --- the failure mode behind the branch-in-base bug of A3, whose
answer flipped between $\mathrm{dps}=15$ and $25$ because $e^{i\pi}=-1+\varepsilon i$ with the
sign of $\varepsilon$ a rounding artifact.

\paragraph{G3 --- quadrature (\code{g3\_quadrature}).}
Path integrals under refinement of the panel count, and \code{res\_circle} under a change of
radius. This is the knob that made the \code{lem:wall} checks miss tolerance while agreeing to
eight digits.

\section*{Provenance}

Over the sessions that produced this suite it found one defect in published work ---
\code{LemRegX3} of arXiv:1806.09874, corrected in \code{fix\_previous\_1806\_file/} --- and none
in \code{finite\_contour\_cocycles\_short.tex}. Everything else it turned up was a defect in the
tests themselves: ill-conditioned coefficient extraction, a branch folded into a base, a
relative-error floor, a winding orientation, an under-resolved residue, a sampled limit, a
$q$-series evaluated outside its disc of convergence, and a patch that silently failed to
apply. That ratio is the argument for keeping it: the machinery under test survived, and the
apparatus built to doubt it did not.

\end{document}
